% Generated draft: requires mathematical and source review.
\documentclass[11pt]{article}
\usepackage[margin=25mm]{geometry}
\usepackage{fontspec}
\setmainfont{Latin Modern Roman}
\usepackage{amsmath,amssymb,mathtools,unicode-math}
\setmathfont{Latin Modern Math}
\usepackage{enumitem,needspace,xurl,hyperref}
\setlength{\parindent}{0pt}
\setlength{\parskip}{5pt}
\setlength{\emergencystretch}{3em}
\newcommand{\N}{\mathbb{N}}
\newcommand{\Z}{\mathbb{Z}}
\newcommand{\Q}{\mathbb{Q}}
\newcommand{\R}{\mathbb{R}}
\newcommand{\C}{\mathbb{C}}
\newcommand{\F}{\mathbb{F}}
\newcommand{\A}{\mathbb{A}}
\newcommand{\PP}{\mathbb{P}}
\newcommand{\E}{\mathbb{E}}
\newcommand{\eps}{\varepsilon}
\newcommand{\dd}{\,\mathrm{d}}
\begin{document}
\section*{UnsolvedMath ID: 30001740}
\textbf{Source problem code:} catalogue entry
\subsection{Definitions and mathematical statement}
The supplied record concerns a spherical homogeneous space \(X=H\backslash G\), but does not specify the ground field, the class of groups, the meaning of sphericity, or the measure defining \(L^2(X)\). Consequently, its intended generality cannot be fixed from this record alone. In the usual algebraic terminology, sphericity means that a Borel subgroup of a reductive algebraic group has an open orbit on the homogeneous space, with the field and geometric interpretation still requiring specification.

The following gives a conditional analytic formulation, rather than asserting hypotheses absent from the source. Suppose that \(G\) is a second-countable locally compact group of type I, that \(H\subseteq G\) is closed, and that \(X=H\backslash G\) has a specified nonzero \(G\)-invariant Radon measure \(\mu_X\). Let \(\widehat G\) be the set of unitary equivalence classes of irreducible unitary representations of \(G\), equipped with the Fell topology. The unitary right regular representation on \(L^2(X,\mu_X)\) is
\[
 (R_X(g)f)(Hx)=f(Hxg).
\]
Its Plancherel decomposition has the form
\[
 L^2(X,\mu_X)\simeq
 \int_{\widehat G}^{\oplus}
 \mathcal H_\pi\otimes\mathcal M_X(\pi)\,\mathrm d\nu_X(\pi),
\]
where \(\mathcal H_\pi\) carries \(\pi\), \(\mathcal M_X(\pi)\) is a multiplicity space, and the spectral measure class is taken on the part where the multiplicity space is nonzero. Define the Fell support by
\[
 \operatorname{Supp}_{\mathrm{Fell}}(R_X)
 =\{[\pi]\in\widehat G:\pi\prec R_X\},
\]
where \(\pi\prec R_X\) denotes weak containment. Explicitly, after fixing a left Haar measure on \(G\), weak containment means
\[
 \|\pi(a)\|\leq\|R_X(a)\|
 \quad\text{for every }a\in C_c(G),
 \qquad
 \pi(a)=\int_G a(g)\pi(g)\,\mathrm dg,
\]
with the analogous definition for \(R_X(a)\).

\textbf{Question.} For the spherical homogeneous spaces intended by the source, characterize precisely
\[
 \operatorname{Supp}_{\mathrm{Fell}}(R_X).
\]
The source also asks which irreducible representations occur in the Plancherel decomposition. This wording needs qualification: a measurable spectral decomposition is specified only up to null sets, whereas Fell support is a closed set and can contain representations without individual positive spectral mass. Thus a classification of almost-everywhere spectral constituents and a characterization of Fell support must be distinguished. If the intended target includes the former, its measure-theoretic meaning must also be specified. The source's claimed equivalence does not resolve this distinction.
\subsection{Short English statement}
For a spherical quotient \(H\backslash G\), determine exactly which irreducible unitary representations belong to the Fell support of the representation on \(L^2(H\backslash G)\). The source leaves the field, group hypotheses, measure, and precise meaning of Plancherel occurrence unspecified; the displayed analytic framework is conditional and requires human review.
\subsection{Sources}
\begin{itemize}
\item Supplied local UnsolvedMath registry record: permanent ID \(30001740\), title \emph{Plancherel Support of Spherical Homogeneous Spaces}, exact source problem code \texttt{catalogue entry}; collection: \emph{Oberwolfach Reports Open Problems}. The record assigns the year \(2011\) and supplies no individual proposer.
\item Supplied report citation: \emph{Automorphic Forms: New Directions} (2011), \emph{Oberwolfach Reports}, pp.~719--779, DOI \texttt{10.4171/owr/2011/14}. Supplied source URL and literature-source link, labeled ``original OWR report'': \url{https://doi.org/10.4171/owr/2011/14}. No problem-specific page or named contributor is supplied.
\item Literature assessment dated exactly \texttt{2026-08-21}: the supplied curation labels the problem open and describes it as a related unresolved question concerning the precise Plancherel support of \(X\)-distinguished representations. This is supplied dated curation, not independently verified; its time-sensitive status and cited claims require independent verification.
\item The registry labels the statement \texttt{corrected\_verified} and reports that its mathematical content was checked against the official report, with notation, context, and question form normalized. This is a supplied verification claim, not an independent verification performed for this draft.
\end{itemize}
\end{document}
